\section{Runtime Checks for Double-Double Arithmetic}
\label{sec:overview}

\lstdefinestyle{sectiontwoc}{
  language=C,
  basicstyle=\ttfamily\normalsize,
  columns=fullflexible,
  keepspaces=true,
  numbers=left,
  numberstyle=\scriptsize\color{gray},
  numbersep=6pt,
  breaklines=true,
  breakatwhitespace=true,
  frame=single,
  framesep=3pt,
  xleftmargin=1.5em,
  framexleftmargin=1.2em,
  framexrightmargin=3pt,
  aboveskip=0pt,
  belowskip=0pt,
  showstringspaces=false
}

This section uses division of a double-double number by a binary64 value to
illustrate a library interface.  We present its two public wrappers, relate
their C data type to the mathematical definition of a double word, and explain
their shared kernel and runtime checks.  We then state their common contract
and summarize the remaining library operations.

\subsection{Runtime Checks for Dividing a Double-Double by a Binary64 Value}
\label{sec:overview-division-wrappers}

The C library provides two public wrappers for dividing a double-double
number by a binary64 value.  Both call the same arithmetic kernel and use the
\texttt{struct dword} type shown in \Cref{fig:dword-struct} to store their
double-double operands and results.

\begin{figure}[t]
  \begin{lstlisting}[style=sectiontwoc]
struct dword {
    double s;
    double t;
};
\end{lstlisting}
  \caption{The C structure used to store a double-double number.  The fields
  contain its two binary64 components.}
  \Description{A C structure containing two binary64 fields named s and t.}
  \label{fig:dword-struct}
\end{figure}

The default public function is the weak-precondition wrapper shown in
\Cref{fig:weak-precondition-division-wrapper}.  The argument \texttt{x}
points to the stored dividend, \texttt{y} is the binary64 divisor, and
\texttt{z} points to storage for the quotient.  Before invoking the kernel,
the wrapper calls
\texttt{dw\_div\_fp\_\allowbreak weak\_precond\_\allowbreak check}.  After the
kernel writes the quotient to \texttt{z}, it calls
\texttt{isfinite\_dword} to check the computed pair.  If either runtime check
fails, the function exits.

\begin{figure}[t]
  \begin{lstlisting}[style=sectiontwoc]
void dw_div_fp(
    struct dword *z, struct dword *x, double y)
{
    if (!dw_div_fp_weak_precond_check(x, y))
        exit(EXIT_FAILURE);

    dw_div_fp_kernel(z, x, y);

    if (!isfinite_dword(z))
        exit(EXIT_FAILURE);
}
\end{lstlisting}
  \caption{The weak-precondition division wrapper.  It checks the operands
  before invoking the kernel and checks result finiteness afterward.}
  \Description{A C function that checks its operands, computes a
  double-double quotient, and checks both result components for finiteness.}
  \label{fig:weak-precondition-division-wrapper}
\end{figure}

Some callers may want the runtime checks to establish result finiteness before
the arithmetic kernel is executed.  The strong-precondition wrapper in
\Cref{fig:strong-precondition-division-wrapper} calls
\texttt{dw\_div\_fp\_\allowbreak strong\_precond\_\allowbreak check} for this
purpose.  If the check
fails, the function exits without invoking the kernel.  If it succeeds, the
kernel writes the quotient directly to \texttt{z}; no result check is needed
afterward.

\begin{figure}[t]
  \begin{lstlisting}[style=sectiontwoc]
void dw_div_fp_strong(
    struct dword *z, struct dword *x, double y)
{
    if (!dw_div_fp_strong_precond_check(x, y))
        exit(EXIT_FAILURE);

    dw_div_fp_kernel(z, x, y);
}
\end{lstlisting}
  \caption{The strong-precondition division wrapper.  Its runtime check
  establishes the required conditions, including result finiteness, before
  the kernel is invoked.}
  \Description{A C function that exits when its strong precondition check
  fails and otherwise calls the division kernel with the caller-provided
  output structure.}
  \label{fig:strong-precondition-division-wrapper}
\end{figure}

The following subsections explain the stored representation, the shared
kernel, and the runtime checks used by these two functions.

\subsection{The Double-Double Data Type}
\label{sec:overview-values}

% Connect struct dword, which is specific to the binary64 C library, to the
% general definition of a double word in Section 1.  Keep the distinction
% explicit: the C fields store the binary64 components; the cited mathematical
% relation determines whether their unevaluated sum is a valid double word.
% End by motivating the C kernel that computes the returned pair.  Do not
% promise a displayed generic algorithm or functional model in Section 2.

\Cref{def:double-word} defines a double word over a generic underlying
floating-point format $\mathbb{F}$.  The C library instantiates $\mathbb{F}$
as binary64 and uses round-to-nearest, ties-to-even.  A double word over
binary64 is a double-double number.  The fields \texttt{s} and \texttt{t} in
\Cref{fig:dword-struct} store its high and low words, respectively.

The C type does not itself enforce the double-word relation.  Because the
numerical guarantee assumes that the dividend is a valid double word, both
wrappers check this property before calling the arithmetic kernel.
\Cref{ex:struct-dword} illustrates how the stored components represent a value
and how this relation is checked.

\begin{example}[Checking a \texttt{struct dword}]
\label[example]{ex:struct-dword}
The following object stores the binary64 values $1$ and $2^{-54}$:

\begin{lstlisting}[language=C]
struct dword x = {
    .s = 1.0,
    .t = 0x1p-54
};
\end{lstlisting}

The value represented by this object is the unevaluated real sum
$1+2^{-54}$.  Under binary64 round-to-nearest, ties-to-even,
$\operatorname{RN}(1+2^{-54})=1$, so the stored pair is a valid double word according to
\Cref{def:double-word}.  Its represented value is not merely the value stored
in \texttt{x.s}; the field \texttt{x.t} retains the correction $2^{-54}$.

The runtime checks used by the wrappers first require both fields
to be finite and then evaluate
\begin{lstlisting}[language=C,numbers=none]
x->s + x->t == x->s
\end{lstlisting}
using binary64 arithmetic.  This comparison succeeds for the object above
because the computed sum rounds to \texttt{x.s}.  Soundness of the check
shows that a successful comparison, together with the finiteness tests,
establishes the mathematical relation in
\Cref{eq:dw-representation}.  In contrast, replacing \texttt{0x1p-54} with
\texttt{0x1p-52} makes the computed sum different from \texttt{x.s}, so the
check fails for the pair.
\end{example}

This test establishes one of the input preconditions required for the
numerical guarantee of the division kernel.  Division has additional
preconditions involving the range of the operands and computations performed
within the kernel.  Their purpose is best understood from the implementation,
so we present the shared C kernel before examining the complete checks
performed by the two wrappers.

\subsection{The Division Kernel}
\label{sec:overview-division-kernel}

The weak-precondition wrapper in
\Cref{fig:weak-precondition-division-wrapper} and the strong-precondition
wrapper in \Cref{fig:strong-precondition-division-wrapper} call the same
arithmetic kernel.
\Cref{fig:dw-div-fp-kernel} shows \texttt{dw\_div\_fp\_kernel}, the C
implementation of the division method given as Algorithm~15 by
\citet{Joldes2017}.

\begin{figure}[t]
  \begin{lstlisting}[style=sectiontwoc]
void dw_div_fp_kernel(
    struct dword *z, struct dword *x, double y)
{
    double th;
    struct dword p;
    double dh;
    double dt;
    double d;
    double tl;

    th = x->s / y;
    fast_2mult(&p, th, y);
    dh = x->s - p.s;
    dt = dh - p.t;
    d = dt + x->t;
    tl = d / y;
    fast_2sum(z, th, tl);
}
\end{lstlisting}
  \caption{The binary64 C kernel for dividing a double-double number by a
  binary64 value.}
  \Description{A C function that computes a leading quotient, forms a
  correction from the reconstructed product and the low word of the dividend,
  and combines the leading quotient and correction using FastTwoSum.}
  \label{fig:dw-div-fp-kernel}
\end{figure}

The kernel first computes the leading quotient \texttt{th} from the high word
of the dividend.  The call to \texttt{fast\_2mult} multiplies this quotient by
the divisor and returns two components used to reconstruct the product.  The
computations of \texttt{dh}, \texttt{dt}, and \texttt{d} use this product and
the low word of the dividend to form a residual.  Dividing the residual by
\texttt{y} produces the quotient correction \texttt{tl}.  Finally,
\texttt{fast\_2sum} combines the leading quotient and correction into the pair
stored in \texttt{z}.

\Cref{sec:numerics} presents \texttt{DWDivFP}, the generic Rocq functional
model that specifies this computation, and proves its numerical guarantees
using Flocq.  As described in \Cref{sec:kernel-refinement}, a VST refinement
proof establishes that \texttt{dw\_div\_fp\_kernel} implements the binary64
instance of \texttt{DWDivFP}.

The refinement result identifies the pair computed by C, but the accuracy and
valid-double-word guarantees require additional mathematical conditions.  The
two wrappers establish these conditions differently.  The following
subsections present their executable C checks, beginning with the
weak-precondition wrapper in
\Cref{fig:weak-precondition-division-wrapper}, which checks the computed
result for finiteness.

\subsection{Checking Finiteness after Division}
\label{sec:overview-post-check}

The weak-precondition wrapper uses the helper functions in
\Cref{fig:finite-checks} to test the operands and the pair computed by the
kernel.  Under strict IEEE~754 semantics, $x-x$ is zero when $x$ is finite
and NaN when $x$ is an infinity or NaN.  Because comparison with NaN returns
false, \texttt{isfinite\_double} returns $1$ exactly when its argument is
finite.  The function \texttt{isfinite\_dword} applies this test to both
stored components.  The complete input check is shown in
\Cref{fig:dw-div-fp-input-check}.

\begin{figure}[t]
  \begin{lstlisting}[style=sectiontwoc]
int isfinite_double(double x)
{
    return (x - x) == 0.0;
}

int isfinite_dword(
    const struct dword *x)
{
    return isfinite_double(x->s) &&
           isfinite_double(x->t);
}
\end{lstlisting}
  \caption{The C functions used to check a binary64 value and both components
  of a \texttt{struct dword} for finiteness.}
  \Description{Two C functions.  The first tests whether a binary64 value is
  finite by subtracting it from itself and comparing the result with zero.
  The second applies this test to both fields of a double-double value.}
  \label{fig:finite-checks}
\end{figure}

\begin{figure}[t]
  \begin{lstlisting}[style=sectiontwoc]
int dw_div_fp_weak_precond_check(
    struct dword *x, double y)
{
    int fin_xh = isfinite_double(x->s);
    int fin_xl = isfinite_double(x->t);
    int fin_y  = isfinite_double(y);

    int dwx = fin_xh && fin_xl &&
              (x->s + x->t == x->s);

    double q_plo = ldexp(1.0, -968);
    int xh_lo = (x->s == 0.0) ||
                (q_plo <= fabs(x->s));

    return dwx && fin_y && xh_lo;
}
\end{lstlisting}
  \caption{The input check used by the weak-precondition division wrapper.
  It tests operand finiteness, dividend validity, and the lower bound required
  by the kernel.}
  \Description{A C function that checks the finiteness of the operands, the
  double-word relation for the dividend, and either a zero high word or a
  lower bound on its magnitude.}
  \label{fig:dw-div-fp-input-check}
\end{figure}

The clause \texttt{dwx} combines finiteness tests for the two stored
components with the comparison illustrated in \Cref{ex:struct-dword}.
The check also tests the divisor for finiteness.  These finiteness tests provide
early rejection for infinities and NaNs; they are not separate hypotheses of
the numerical theorem.

The remaining condition permits $x_h=0$ or requires
$|x_h|\geq2^{-968}$.  In the first case, the double-word relation
checked by \texttt{dwx} also implies $x_\ell=0$, so the dividend is exactly
zero.  For a nonzero high word, the lower bound is required by the product
reconstruction performed by \texttt{fast\_2mult} in
\Cref{fig:dw-div-fp-kernel}.  \Cref{sec:division-numerics} derives this
requirement from the numerical analysis.

The input check imposes no upper bound on the operands or leading quotient.
After the kernel executes, \texttt{isfinite\_\allowbreak dword} checks both
components of the computed pair.  It succeeds for normal values, subnormal
values, and zeros, and fails if either component is an infinity or NaN.

\subsection{Guaranteeing Finiteness before Division}
\label{sec:overview-strong-precondition-check}

The strong-precondition wrapper extends the preceding input check with
bounds that establish result finiteness before the kernel executes.  Its
complete check is shown in \Cref{fig:dw-div-fp-check}.  The additional
comparisons bound the high word of the dividend and the rounded leading
quotient.

\begin{figure}[t]
  \begin{lstlisting}[style=sectiontwoc]
int dw_div_fp_strong_precond_check(
    struct dword *x, double y)
{
    int fin_xh = isfinite_double(x->s);
    int fin_xl = isfinite_double(x->t);
    int fin_y  = isfinite_double(y);

    int dwx = fin_xh && fin_xl &&
              (x->s + x->t == x->s);

    double q_omega = ldexp(1.0, 1022);
    double q_plo   = ldexp(1.0, -968);

    int xh_hi = fabs(x->s) <= q_omega;
    int xh_lo = (x->s == 0.0) ||
                (q_plo <= fabs(x->s));

    double q = x->s / y;
    int qt_hi = fabs(q) < q_omega;

    return dwx && fin_y &&
           xh_hi && xh_lo && qt_hi;
}
\end{lstlisting}
  \caption{The input check used by the strong-precondition division wrapper.
  In addition to the tests in \Cref{fig:dw-div-fp-input-check}, it bounds the
  high word and rounded leading quotient.}
  \Description{A C function that checks the operands and requires upper
  bounds on the high word and rounded leading quotient before allowing the
  division kernel to execute.}
  \label{fig:dw-div-fp-check}
\end{figure}

The first additional test requires $|x_h|\leq 2^{1022}$.  The second computes
the same leading quotient as the first line of
\texttt{dw\_div\_fp\_kernel}, $t_h=\operatorname{RN}(x_h/y)$, and requires
$|t_h|<2^{1022}$.  The C variable \texttt{q} stores this rounded
leading quotient.  In particular, the strict comparison fails for a nonfinite
leading quotient.  The two upper bounds are sufficient to show before
execution that the computations in the kernel and both components of its
result remain finite.  No lower bound is imposed on the leading quotient, so
the check permits subnormal values.

These bounds are sufficient but not necessary.  Checks based on them may
therefore fail for operands whose kernel result would be a finite pair.  This
conservative rejection is the cost of establishing finiteness before execution.
\Cref{sec:numerics} derives the bounds and connects them to the generic
numerical conditions.

\subsection{The Division Contract}
\label{sec:overview-division-guarantee}

Although the weak-precondition wrapper in
\Cref{fig:weak-precondition-division-wrapper} and the strong-precondition
wrapper in \Cref{fig:strong-precondition-division-wrapper} establish
result finiteness at different points, they provide the same contract on
normal return.

\begin{theorem}[Normal-Return Guarantee for the Division Wrappers]
\label{thm:guarded-division}
Suppose that \texttt{x} initially stores $(x_h,x_\ell)$ and that a call to
either \texttt{dw\_div\_fp} or
\texttt{dw\_div\_fp\_\allowbreak strong} returns normally
with \texttt{z} storing $(z_h,z_\ell)$.  Writing
$q=(x_h+x_\ell)/y$ for the exact real quotient, the represented result
satisfies
\[
  \left|(z_h+z_\ell)-q\right|
  \leq 3u^2|q|+\eta,
\]
where $u=2^{-53}$ is the unit roundoff and $\eta=2^{-1074}$ is the smallest
positive subnormal value for binary64.  The returned pair is also a valid
double-double according to \Cref{def:double-word}.
\end{theorem}

Callers need not supply numerical preconditions.  Instead, each wrapper
establishes the sufficient conditions required by the underlying
numerical result.  The guarantee is conditional on normal return: if any
check fails, the wrapper follows its nonreturning error path and provides no
numerical result.

\Cref{sec:mechanization} shows how the VST proofs combine soundness of the
corresponding check, refinement of the C kernel, and the numerical theorem to
establish this contract.

\subsection{Operations and Guarantees of the C Library}
\label{sec:overview-library-guarantees}

The two division wrappers illustrate the form of the contract provided for
each double-word operation in the C library.  \Cref{tab:operation-preview}
places division alongside the other available operations and summarizes the
accuracy guarantee provided when a wrapper returns normally.

\begin{table*}[t]
\centering
\caption{Normal-return accuracy guarantees provided by the C library.  Here
$a$ and $b$ are binary64 operands, $x$ and $y$ are the values represented by
double-double operands, $r$ is the exact result, and
$\widehat r=z_h+z_\ell$ is the value represented by the returned pair.  For
binary64, the unit roundoff is $u=2^{-53}$ and the smallest positive
subnormal value is $\eta=2^{-1074}$.  Every returned pair is also a valid
double-double number.}
\Description{The operations provided by the C library, their exact real
results, and their accuracy guarantees when a wrapper returns normally.}
\label{tab:operation-preview}
\small
\begin{tabular}{@{}p{0.49\textwidth}p{0.45\textwidth}@{}}
\toprule
Operation & Normal-return accuracy guarantee \\
\midrule
\multicolumn{2}{@{}l}{\emph{Error-free transformations}} \\[2pt]
\TwoSum{}$(a,b)$ and \FastTwoSum{}$(a,b)$; $r=a+b$
  & $\widehat r=r$ \\
\FastTwoMult{}$(a,b)$; $r=ab$
  & $\widehat r=r$ \\
\addlinespace
\multicolumn{2}{@{}l}{\emph{Double-word operations}} \\[2pt]
Addition $(x,a)$; $r=x+a$
  & $|\widehat r-r|\leq 2u^2|r|$ \\
Addition $(x,y)$; $r=x+y$
  & $|\widehat r-r|
       \leq \dfrac{3u^2}{1-4u}|r|$ \\
Multiplication $(x,a)$; $r=xa$
  & $|\widehat r-r|
       \leq 2u^2|r|+2.51\eta$ \\
Multiplication $(x,y)$; $r=xy$
  & $|\widehat r-r|
       \leq 4u^2|r|+2.5\eta$ \\
Division $(x,a)$; $r=x/a$
  & $|\widehat r-r|
       \leq 3u^2|r|+\eta$ \\
Division $(x,y)$; $r=x/y$
  & $|\widehat r-r|
       \leq 9.8u^2|r|+2.5\eta$ \\
\bottomrule
\end{tabular}
\end{table*}

Each of the six double-word operations in \Cref{tab:operation-preview} is
available through both wrappers.  The two wrappers provide the same
guarantee shown in the table, although their checks can succeed on different
sets of inputs.
Each error-free transformation instead has one wrapper for its
exactness guarantee.

The set of inputs for which all runtime checks succeed depends on which
computations in the corresponding algorithm must remain finite and
how rounding near the lower end of the format affects its result.  Determining
these domains therefore requires a numerical analysis that accounts for the
finite exponent range of IEEE~754 arithmetic.  \Cref{sec:numerics} develops
this analysis and derives sufficient conditions for each error bound and valid
double-word result.
